%% ======================================================================
%% Section 3 "Data" -- drop-in replacement for the empty \section{Data}
%% block in main.tex. Two subsections as requested:
%%   3.1 Synthetic dataset creation  -- full methodology text, ready to use
%%   3.2 Dataset description        -- skeleton only, left for the author
%%       to complete once final counts/figures are available
%% Bracketed placeholders [LIKE THIS] mark numbers that depend on the
%% final production build (still pending as of this draft) and must be
%% filled in once that run completes -- do not guess these before then.
%% ======================================================================

\subsection{Synthetic dataset creation}
\label{subsec:dataset_creation}

The synthetic ground-truth dataset used in this study is derived entirely
from network-side data already collected by the distribution network
operator: a circuit-level interval telemetry table (hereafter the
\emph{source telemetry table}) and an accompanying circuit/site metadata
table. The source telemetry table records 5-minute interval readings of
active power, reactive power (derived from cumulative reactive energy),
voltage, current, power factor, and cumulative imported/exported energy
for every monitored circuit at every site, partitioned by year, month, and
generation/load side. Over the [2024--2025] study period it spans
approximately [X] billion rows across [Y] GB. The metadata table carries,
per circuit, a \emph{circuit type} label distinguishing the whole-site
load-side signal from the whole-site PV-side signal (plus a set of
named, partial sub-load circuits not required for this study), a
\emph{circuit polarity} indicator, a \emph{device identifier} grouping
co-located phases of the same physical meter, and site-level attributes
including two independent PV capacity proxies (a nameplate rating and an
empirical 99th-percentile apparent-power figure) used later for
capacity-normalised comparisons (Section~\ref{subsec:capacity_proxies}).

\subsubsection*{Cohort selection}
Two circuit types anchor the ground truth used throughout this paper: a
load-side circuit type, representing the site's net load as measured at
the load-side conductor, and a PV-side circuit type, representing the
site's PV generation. A candidate site was retained only if it exposed at
least one circuit of each type in the metadata table. Because single- and
three-phase installations are both present in the fleet, and the
metadata table carries no verified circuit-to-phase (A/B/C) label for
either side, sites were further restricted to those with exactly one or
exactly three load-side circuits -- the two circuit counts consistent
with a genuine single- or three-phase installation -- deferring the
remaining, structurally ambiguous sites (multi-device sites and sites
with two or four+ load-side circuits) to future work rather than guessing
at their topology. This restriction yielded [14{,}353] clean candidate
sites out of [15{,}048] sites exposing at least one load-side circuit.

\subsubsection*{Data-quality issues identified and handled}
An initial signal-taxonomy investigation surfaced five issues that, left
uncorrected, would have silently biased the resulting ground truth:

\begin{enumerate}
  \item \textbf{Duplicate circuits.} A material share of sites exposed
    two or more circuit identifiers reporting near-identical (correlation
    $\geq 0.99$) readings, both within the same circuit type (e.g. two
    \texttt{PV} circuit identifiers reporting the same underlying signal)
    and across types. Left unresolved, these would double-count load or
    generation. A per-site duplicate-detection step identifies and drops
    the redundant identifier(s), keeping one representative circuit per
    physical signal.
  \item \textbf{Inactive circuits.} A smaller share of circuits reported
    a near-zero (sub-few-Watt) reading for an entire day, consistent with
    a decommissioned or never-commissioned meter channel still present in
    the metadata. These are detected and dropped independently of the
    duplicate check.
  \item \textbf{A device/meter-model power-reporting artefact.} A subset
    of circuits, strongly associated with a specific meter/device model,
    reports an instantaneous power field that is unreliable at the
    nominal 5-minute cadence: the underlying energy register only
    accumulates in whole Watt-hour steps at a true interval closer to
    five minutes, and reading power directly off it aliases against the
    true accumulation interval. For these circuits, power is instead
    re-derived from the accumulated energy and the circuit's own implied
    accumulation interval; this correction is applied by default but is
    exposed as a configurable option, since it is what makes this
    dataset's active-power series diverge (visibly, since the correction
    produces a step-wise rather than smooth trace) from a naive
    energy-agnostic read of the same source table.
  \item \textbf{Sign convention.} Raw per-circuit power and reactive
    power are not consistently signed with respect to import/export
    across circuits. A per-circuit polarity correction, derived and
    verified against known-generation and known-load circuits, is
    applied uniformly so that positive PV-side power always denotes
    generation and load-side power reflects genuine net (import-minus-
    self-consumed-generation) flow at the load conductor, consistent
    with a physical net meter rather than a pure consumption meter.
  \item \textbf{Night-time reconstruction failures.} For every candidate
    site, load-side and PV-side signals were summed (after polarity
    correction) and checked against an expected near-zero PV contribution
    overnight. A small share of sites showed an implausible non-zero
    night-time reconstructed load, consistent with an undeclared
    behind-the-meter storage circuit or an uncorrected sign issue outside
    the scope of the polarity correction above. Rather than attempt a
    per-site correction, these sites were excluded entirely, along with
    any site whose duplicate-circuit resolution could not be settled
    unambiguously (flagged for manual review) -- a deliberately
    conservative choice, favouring a smaller but trustworthy dataset over
    a larger one requiring untested per-site assumptions.
\end{enumerate}

Because neither the metadata table nor the source telemetry table carries
a verified circuit-to-phase label, one further modelling choice was
necessary wherever PV-side and load-side circuit counts disagree at a
site (e.g. one PV-side circuit serving a three-phase load): PV generation
is allocated evenly across the site's load phases in that case, and
matched one-to-one, by circuit identifier ordering, only where the two
counts agree. This is treated explicitly as a heuristic, not a verified
pairing, and every derived value carries a record of which rule produced
it (Section~\ref{subsec:dataset_description}).

\subsubsection*{Construction pipeline}
The final dataset was built in four stages, applied to the [14{,}353]
clean candidate sites identified above:

\begin{enumerate}
  \item \textbf{Fleet-wide resolution (single representative day).} The
    duplicate, inactivity, and reconstruction checks above were run
    across the full candidate cohort for one representative day, yielding
    a per-circuit keep/drop decision and a per-site exclusion decision
    (manual-review and reconstruction-failure sites dropped entirely).
  \item \textbf{Full-period extraction.} For every circuit surviving
    stage 1, the complete [2024--2025] interval history was extracted
    from the source telemetry table, restricted to the load-side and
    PV-side circuit types only (the named partial sub-load circuits were
    not required for this study).
  \item \textbf{Full-history revalidation.} The single-day checks in
    stage~1 establish circuit and site quality on one day only; a circuit
    or site can still degrade later in the study period (e.g. a device
    failure, a newly commissioned behind-the-meter battery). Stage~3
    therefore repeats the inactivity and reconstruction checks across
    every landed month rather than the single sampled day: a circuit
    found newly inactive in any month is dropped from that point forward;
    a site newly failing the reconstruction check in any month is
    excluded from the dataset entirely.
  \item \textbf{Table construction.} From the fully revalidated circuit
    set, the three deliverable tables described in
    Section~\ref{subsec:dataset_description} are built one landed month
    at a time, with the device/meter-model and sign-convention
    corrections above applied throughout.
\end{enumerate}

\subsubsection*{Final dataset}
The output of this pipeline is three related tables, each built from the
same underlying revalidated circuit-level telemetry but aggregated
differently, so that the signal an algorithm is evaluated \emph{against}
is never derived from the same aggregation as the signal it is given as
\emph{input}:

\begin{itemize}
  \item A \textbf{site-level ground-truth table}, one row per site per
    timestamp, giving load-side and PV-side active power, reactive power,
    and voltage separately, plus their combination into a single
    PV-independent gross load figure and a capacity-normalised PV output
    figure. This is the target a disaggregation algorithm's site-level
    output is graded against.
  \item A \textbf{synthetic smart-meter table}, one row per site, device,
    circuit, and timestamp, restricted to load-side circuits only and
    kept per phase (not pre-summed), using the raw, uncorrected
    per-circuit reading. This is the only table intended as algorithm
    \emph{input} -- it is deliberately built to resemble what a
    distribution network operator's existing AMI infrastructure already
    reports, without exposing the independently metered PV signal an
    algorithm is meant to infer.
  \item A \textbf{per-circuit ground-truth table}, one row per site,
    device, circuit, and timestamp, for every surviving load-side
    \emph{and} PV-side circuit, sign-corrected but otherwise
    unaggregated. This supports per-phase grading and topology-aware
    analysis, and carries the PV-to-phase allocation rule
    (Section~\ref{subsec:dataset_creation}) as an explicit per-site tag
    rather than folding it into a derived numeric value.
\end{itemize}

Across the [2024--2025] period and the sites surviving every stage above,
the final dataset comprises [X] sites, [Y] rows in the site-level table,
and [Z] rows in the synthetic smart-meter table (Section~\ref{subsec:dataset_description}
gives the full breakdown).

\subsection{Dataset description}
\label{subsec:dataset_description}

%% TODO (author): complete this subsection once the final [2024--2025]
%% production build has run. Suggested structure, matching the three
%% tables introduced in Section~\ref{subsec:dataset_creation}:
%%
%%   - Fleet size and attrition: candidate sites -> clean cohort ->
%%     sites surviving into each of the three final tables, with the
%%     reasons for attrition at each stage (duplicate/inactive circuits,
%%     manual-review/reconstruction exclusions, sites with no temporal
%%     overlap between load and PV signals).
%%   - Geographic distribution (e.g. postcode/DNSP-zone choropleth or
%%     scatter, if location data is included/permitted).
%%   - PV and load capacity distribution across the final cohort (e.g.
%%     histogram of nameplate/S_99 capacity; single- vs three-phase
%%     PV/load split; counts of the four load/PV phase-count
%%     combinations, e.g. 1-phase load with 3-phase PV).
%%   - Temporal coverage per site (date range actually landed per site,
%%     given real-world circuit/device churn -- some sites' data ends
%%     before the nominal study end date).
%%   - Table-level summary statistics: row counts, site counts, and
%%     total data volume for each of the three final tables (see the
%%     bracketed placeholders left in Section~\ref{subsec:dataset_creation}).
%%   - A short note on data provenance/reproducibility: the pipeline
%%     above is deterministic and re-runnable end-to-end from the source
%%     telemetry table, with every exclusion decision retained in an
%%     audit trail rather than only reflected in the final row counts.

\hspace{1em} % placeholder paragraph break -- replace with the text above
